% =====================================================================
\section{The big picture in plain words}
% =====================================================================

\subsection{The question the project answers}

Imagine a long thin metal rod. At the start, some parts are hot and some are cold. You cannot watch the whole
rod continuously; instead, a few hundred cheap, noisy thermometer readings were taken at random places and
random times. Each reading is just three numbers:
\[
  (x,\; t,\; u) \;=\; (\text{position},\; \text{time},\; \text{measured temperature}).
\]
The question is: \textbf{what rule of nature produced these numbers?} Not ``draw a curve through the points'',
but ``find the \emph{law}'', the equation that says how $u$ changes in time at every point, so that you could
predict what happens at times and places nobody measured.

That is what this project calls \emph{scientific discovery}: given sparse, noisy observations of a field
$u(x,t)$, and \emph{not} the governing equation, recover the equation, estimate its unknown constants, reject
explanations that do not fit, and use the winner to predict outside the observed window.

\begin{intuition}[: a detective, not a curve-fitter]
A curve-fitter answers ``what numbers come next?''. A detective answers ``who did it, and how?''. This project
is built like a detective agency:
\begin{itemize}
  \item a \textbf{librarian} (RAG) fetches relevant textbook pages about how fields usually behave;
  \item an \textbf{archivist} (knowledge graph) keeps a structured card index of which mechanisms use which
        equations and parameters;
  \item a \textbf{brainstormer} (the LLM) suggests suspects: candidate equations;
  \item a \textbf{forensic lab} (physics-informed neural networks) tests each suspect against the evidence and
        estimates its unknown constants;
  \item a \textbf{judge} (model selection) picks the simplest suspect that fully explains the evidence;
  \item and a final \textbf{prediction test} checks the verdict on data that nobody used for any decision.
\end{itemize}
Crucially, the brainstormer is never allowed to decide the verdict. Only measured evidence decides.
\end{intuition}

\subsection{What the answer turned out to be}

The data was generated synthetically, so the true answer is known (but hidden from every decision):
\[
  \frac{\partial u}{\partial t} \;=\; D\,\frac{\partial^2 u}{\partial x^2}, \qquad D = 0.1 .
\]
This is the \emph{diffusion} (or heat) equation. In words: \emph{each point drifts towards the average of its
neighbours, at a speed set by $D$}. The project recovered this equation from 400 noisy points, estimated
$\hat D = 0.10002$ (0.02\% error), rejected the wrong candidates, and predicted the field beyond the observed
time window with an error of $2.5\times10^{-4}$. Phase~7 then repeated the whole pipeline on a different,
unseen system (advection--diffusion) and recovered both of its constants within 0.4\%.

\subsection{How the project was built: eight phases}

The project grew in phases, each adding one capability and each \emph{reporting its failures honestly}:

\begin{center}
\small
\begin{tabular}{@{}p{1.4cm}p{4.3cm}p{8.9cm}@{}}
\toprule
\textbf{Phase} & \textbf{Adds} & \textbf{One-line takeaway} \\ \midrule
1 & Synthetic data & Simulate diffusion, then sample 400 noisy points. A single-sine start turned out to be a trap. \\
2 & Equation discovery (SINDy) & Regression on numerical derivatives finds the law on clean data, fails on noisy data. \\
3 & PINNs & A neural network that must obey a candidate equation; it validates hypotheses and estimates $D$. \\
4 & Hypothesis generation (LLM) & An LLM proposes equations as structured JSON; a compiler turns any equation into a PINN. \\
5 & RAG + knowledge graph & Ground the LLM in retrieved textbook evidence, organised as a graph with provenance. \\
6 & Physics-constrained search & The full pipeline: pool $\to$ filters $\to$ PINN per candidate $\to$ selection $\to$ OOD. \\
7 & Evaluation & 12 benchmark datasets, 95 questions, 20 adversarial attacks, ablations, a second system. \\
8 & Final report & Audit of every claim: what was demonstrated and what was not. \\
+ & Real LLM (Gemini) & The mock LLM was replaced by Google Gemini for a real run (Section~\ref{sec:gemini}). \\
\bottomrule
\end{tabular}
\end{center}

\subsection{How to read this guide}

\begin{itemize}
  \item \textbf{Section~\ref{sec:background}} teaches every concept from zero. If a word in a later section is
        unfamiliar, it is explained there.
  \item \textbf{Section~\ref{sec:architecture}} shows the whole machine on one page, plus a map of how the code
        packages depend on each other.
  \item \textbf{Sections~\ref{sec:p1}--\ref{sec:p7}} walk through the phases. Each has a function-level flowchart.
        Boxes are coloured by the package the function lives in (legend under each chart). Solid arrows mean
        ``calls / passes data to''; dashed arrows mean ``feeds back'' or ``used later''.
  \item \textbf{Section~\ref{sec:gemini}} covers the real Gemini run and compares it with the mock LLM.
  \item \textbf{Appendix~\ref{app:functions}} is a complete function reference you can look things up in.
\end{itemize}
